\documentclass[twocolumn]{aastex631}

\usepackage{amsmath}
\usepackage{hyperref}

\newcommand{\Gravitas}{\textsc{Gravitas}}
\newcommand{\Version}{1.0.0}
\newcommand{\ScenarioCount}{59}
\newcommand{\InvestigationCount}{22}
\newcommand{\InvestigationSteps}{631}
\newcommand{\ReleaseDOI}{ZENODO-DOI-PLACEHOLDER}
\newcommand{\ZenodoRecord}{ZENODO-RECORD-URL-PLACEHOLDER}
% Full source snapshot: a4e6ab3a5ecb69290742c7197ed545737136182a
\newcommand{\CodeSnapshot}{a4e6ab3}

% This source intentionally compiles before the four final screenshots are
% present. Add each PNG named below to paper/figures/ and the placeholder panel
% will automatically be replaced by the image.
\newcommand{\figureplaceholder}[2]{%
  \begin{center}
  \fbox{\begin{minipage}[c][#1][c]{0.92\linewidth}
  \centering\small\textbf{Final screenshot to insert}\\[4pt]#2
  \end{minipage}}
  \end{center}
}

\shorttitle{Gravitas}
\shortauthors{Ziegler}

\begin{document}

\title{Gravitas: A Browser-Based Astrophysics Laboratory for Prediction,
Measurement, and Discovery}

\author{Carl Ziegler}
\affiliation{Department of Physics, Engineering and Astronomy, Stephen F. Austin
State University, Nacogdoches, TX 75962, USA}
\email{Carl.Ziegler@sfasu.edu}

% Replace the ORCID placeholder above if an ORCID is desired in the arXiv copy.

\begin{abstract}
Astronomy students are often asked to reason about systems that cannot be
handled in a laboratory, evolve over inaccessible timescales, or reveal
themselves only through indirect measurements. \Gravitas\ is a free,
browser-based astrophysics environment designed to turn those limitations into
opportunities for experimentation. Its planned version \Version\ release
combines a configurable gravitational sandbox with \ScenarioCount\ ready-to-run
scenarios and \InvestigationCount\ guided investigations. Students can build and
alter systems, change reference frames, measure distances, angles, periods,
energies, forces, and orbital elements, and observe simulated systems through
transit, radial-velocity, astrometric, rotation-curve, stellar-evolution, and
gravitational-wave displays. Guided activities use a
predict--test--measure--revise--explain cycle, retain the student's evidence, and
export a report for submission. Instructor guides, answer keys, short classroom
routes, lecture and embed modes, reproducible share links, Spanish localization,
keyboard support, and offline operation lower the practical barrier to adopting
the simulation in introductory astronomy and physics courses. This article
describes the learning experiences available in \Gravitas, illustrates several
ways it can be used from a brief demonstration to a full laboratory exercise,
and states the scientific scope and limitations of the models. The aim is not to
replace quantitative calculation or observation, but to give students a place
where they can make a prediction, change one thing, measure what follows, and
build an explanation from their own evidence.
\end{abstract}

\keywords{Astronomy education (2165) --- Astronomy software (1855) ---
Computational methods (1965) --- Open source software (1866) ---
Physics education --- Interactive learning environments}

\section{Introduction}
\label{sec:introduction}

Astronomy is unusually rich in phenomena that students can see but cannot
directly manipulate. They cannot move a planet closer to its star, turn off a
galaxy's dark matter, replay a close stellar encounter, or wait through a
main-sequence lifetime. Even familiar ideas such as an elliptical orbit or a
transit light curve require students to coordinate several representations at
once: a moving physical system, a graph, a mathematical relationship, and the
viewpoint of an observer.

Interactive simulations can make such systems available for inquiry when the
simulation is paired with purposeful prompts, measurements, and opportunities
to explain the result \citep{finkelstein2005,wieman2008,rutten2012}. Active
engagement is not created by animation alone; students need a question to
answer and evidence with which to answer it \citep{hake1998,freeman2014}.
\Gravitas\ was built around that distinction. It can be opened as a free-form
sandbox, but it also contains a complete instructional layer that asks students
to predict an outcome before running a system, measure the result with
task-specific instruments, compare it with their prediction, and explain what
changed.

The intended audience is broad: instructors teaching introductory astronomy or
physics, students exploring independently, and presenters who want a
reproducible live demonstration. A visitor needs no account, installation, or
specialized hardware. The simulation runs in a modern web browser at
\url{https://gravitas-sim.online}; the source is openly available under the MIT
License. A complete state can be encoded in a URL, allowing an instructor to
distribute the same initial conditions to an entire class and allowing a student
to return a modified system as part of an answer.

This paper presents \Gravitas\ as a teaching environment rather than as a
software architecture. Section~\ref{sec:experience} introduces the student
experience. Sections~\ref{sec:investigations} and \ref{sec:tools} describe the
guided investigations and measurement tools. Section~\ref{sec:examples} follows
three representative learning experiences, while
Section~\ref{sec:instructors} outlines practical classroom use.
Section~\ref{sec:scope} states what the models can and cannot support. The paper
describes the release candidate represented by source commit
\texttt{\CodeSnapshot}; counts and screenshots should be refreshed from the
tagged \Version\ release before submission.

\section{From First Click to Experiment}
\label{sec:experience}

\Gravitas\ opens directly into an interactive system. The central canvas shows
the evolving scene, while controls expose bodies, scenario settings, time,
viewpoint, and measurement tools. A learner can start from the Solar System,
TRAPPIST-1, a binary star, a spacecraft transfer, a galactic rotation curve, or
a black-hole encounter, then change the masses, positions, velocities, and
physical options that define the experiment. The current gallery contains
\ScenarioCount\ systems spanning orbital motion, exoplanets, stars, galaxies,
compact objects, and playful counterfactual cases. Concept tags let an
instructor find a system that matches the topic of a particular class meeting.

The interface is designed to move naturally between qualitative observation and
quantitative evidence. Trails reveal orbital shape and motion. Velocity and
acceleration arrows expose the changing geometry of force and motion. A
selection inspector reports mass, radius, speed, orbital elements, and energy
for an individual body. A ruler, protractor, stopwatch, plots, and event-aware
measurements let students attach numbers to what they see. The elapsed
simulation time and spatial scale remain visible, and exported screenshots
carry the scenario name, clock, and scale so that an image remains interpretable
after it leaves the application.

\begin{figure*}
\centering
\IfFileExists{figures/overview.png}
  {\includegraphics[width=0.96\textwidth]{figures/overview.png}}
  {\figureplaceholder{2.35in}{A wide overview of the \Gravitas\ workspace using
  the Kepler's Second Law scenario: simulation canvas, object controls,
  selected-body inspector, ruler or stopwatch, scale bar, elapsed time, and
  velocity/acceleration arrows visible at once.}}
\caption{The \Gravitas\ workspace is both a visual sandbox and a measurement
environment. The final image should show a scientifically legible scenario
rather than a decorative collision: the orbit, clock, scale, vectors, and at
least one measurement tool should all be visible.}
\label{fig:overview}
\end{figure*}

A scenario is not a locked animation. Students can pause, step, rewind by
restoring a captured state, or change the rate at which simulated time passes.
They can follow a body with the camera or re-express the entire recorded scene
in that body's reference frame. The latter distinction matters: changing to
Earth's frame redraws the prior trails and causes Mars to trace the familiar
retrograde loop, while the underlying heliocentric dynamics remain unchanged.
The student therefore sees reference frame as a choice of description rather
than a change in physical law.

The same emphasis on reproducibility appears throughout the application.
Shareable links preserve the initial configuration. Seeded scenarios reproduce
the same controlled trial. An A/B experiment bench captures a state, records a
baseline, restores the identical state, and then records a second run after one
variable is changed. The resulting comparison names the changed parameter and
can be exported with provenance. These features support a simple but important
classroom habit: change one thing, preserve everything else, and decide whether
the evidence supports the claimed cause.

\section{Guided Investigations}
\label{sec:investigations}

The \InvestigationCount\ guided investigations contain
\InvestigationSteps\ student-facing steps. They range from short,
15--25 minute spaceflight activities to 80--100 minute stellar-evolution
sequences. All are aimed primarily at introductory astronomy, with one
gravitational-wave investigation explicitly requiring no prior physics.

Each investigation follows a recurring learning cycle:

\begin{enumerate}
\item \textbf{Predict.} The learner commits to an expected outcome before the
simulation reveals it. Predictions are recorded but are not graded for
correctness.
\item \textbf{Test.} A prepared scene loads with the controls and instruments
needed for the question.
\item \textbf{Measure.} The learner records values from the running model rather
than copying numbers from explanatory text.
\item \textbf{Revise.} The learner compares the result with the original
prediction and, where appropriate, changes a variable or repeats a controlled
run.
\item \textbf{Explain.} A short response, calculation, graph, or comparison asks
the learner to connect the evidence to the physical idea.
\end{enumerate}

This structure is intended to limit passive watching. The prediction gives the
result something to challenge; the measurement makes the conclusion auditable;
and the final explanation requires the student to say why the result occurred.
Progress is stored locally, and the finished evidence notebook can be exported
as a PDF for an existing learning-management system. No student account or
external data service is required.

\begin{figure}
\centering
\IfFileExists{figures/investigation-workflow.png}
  {\includegraphics[width=\columnwidth]{figures/investigation-workflow.png}}
  {\figureplaceholder{2.55in}{A vertically arranged composite from one
  investigation showing (1) a prediction prompt, (2) an attached live
  measurement, and (3) the evidence/report summary. Use one continuous example,
  preferably Kepler's Laws or Finding Planets by Their Shadows.}}
\caption{A guided investigation turns interaction into a record of reasoning.
Students predict before the run, attach measurements from the live scene, and
finish with an exportable evidence-based explanation.}
\label{fig:workflow}
\end{figure}

Table~\ref{tab:investigations} summarizes the breadth of the catalogue. The
investigations are not merely tours of features: each is organized around a
claim that students can settle with a measurement. For example, the resonance
activity first demonstrates that a near-integer period ratio is not sufficient
evidence for resonance, then asks students to distinguish libration from
circulation of the resonant angle. The binary-planet activity maps the boundary
between surviving and disrupted orbits. The gravitational-wave sequence asks
students to distinguish what is physically calculated, what is an illustrative
visualization, and what an observer measures.

\begin{table*}[t]
\caption{Guided investigation catalogue}
\label{tab:investigations}
\centering
\small
\renewcommand{\arraystretch}{1.12}
\begin{tabular}{p{0.17\textwidth}p{0.31\textwidth}p{0.45\textwidth}}
\hline\hline
\textbf{Theme} & \textbf{Investigations} & \textbf{Representative student work}\\
\hline
Orbits and gravity &
Kepler's Laws; Why Mars Goes Backwards; Bound, Unbound and Escape; Weighing the
Stars; Tides &
Locate the focus of an ellipse; compare equal swept areas; transform reference
frames; classify trajectories by energy; infer stellar mass from a binary
orbit; measure differential gravity.\\
Exoplanets and observing &
Finding Planets by Their Shadows; Finding Planets by Their Tug; The Goldilocks
Question; Can You Detect This Planet?; Design the Schedule; Planets in Binary
Stars &
Recover radius and period from a transit; infer mass from radial velocity;
combine mass and radius; explore inclination and dilution; design observations;
map stable regions around one or two stars.\\
Orbital dynamics and spaceflight &
The Butterfly Effect in Space; When Orbits Lock; Where Does a Gravity Assist Get
Its Speed?; Getting There From Here; Where Can It Get To? &
Measure divergence of nearby starts; test resonant angles; compare a flyby in
two frames; calculate and fly a Hohmann transfer; use the Jacobi constant and
zero-velocity regions.\\
Compact objects and gravitational waves &
Black Holes by the Numbers; What Is a Gravitational Wave?; Listening to
Spacetime &
Test how black-hole properties scale with mass; connect orbital motion to a
travelling strain; measure a chirp and compare the model with GW150914
\citep{abbott2016}.\\
Stars, galaxies, and dark matter &
A Universe of Stars; Lives of Stars; The Missing Mass &
Separate mass, radius, temperature, and luminosity on the H--R diagram; follow
published stellar tracks \citep{dotter2016,choi2016}; fit a rotation curve and
compare luminous and dynamical mass.\\
\hline
\end{tabular}
\vspace{3pt}

\parbox{0.95\textwidth}{\footnotesize\textit{Note.} Durations and prerequisites
are supplied in the application. The catalogue can be filtered by concept, and
individual investigations may be assigned through reproducible links.}
\end{table*}

\section{Tools That Turn Motion into Evidence}
\label{sec:tools}

The tools in \Gravitas\ are selected for questions commonly asked in
introductory astronomy. They are available in the sandbox and are also opened
automatically when a guided step requires them.

\begin{table*}[t]
\caption{Student tools and the reasoning they support}
\label{tab:tools}
\centering
\small
\renewcommand{\arraystretch}{1.12}
\begin{tabular}{p{0.21\textwidth}p{0.34\textwidth}p{0.38\textwidth}}
\hline\hline
\textbf{Tool} & \textbf{What the student records} & \textbf{Example question}\\
\hline
Ruler, protractor, and scale bar &
Distances in astronomical units or kilometres; angles between locations or
vectors &
Where is the focus? How far does a planet move in a fixed interval? What angle
separates a Trojan asteroid from Jupiter?\\
Simulation-time stopwatch &
Elapsed time, including automatic latching at successive periapsis passages &
How can a period be measured without relying on reaction time?\\
Object inspector &
Mass, radius, speed, orbital elements, and kinetic, potential, and total energy &
Is the trajectory bound? Which body is the orbit actually referenced to?\\
Velocity, acceleration, and source arrows &
The selected body's velocity, total acceleration, and contribution from each
gravitating source &
Why is an orbit curved? Why are velocity and force not generally parallel?\\
Reference-frame controls &
The same positions and recorded trails relative to the barycenter or a selected
body &
Why does Mars reverse direction in an Earth-centered view?\\
A/B experiment bench &
Aligned baseline and comparison time series, changed parameter, CSV, and
provenance metadata &
Does changing the starting distance, rather than the random seed or timestep,
cause the different outcome?\\
Observing panels &
Transit flux, radial velocity, astrometric displacement, viewing geometry, and
sampled observations &
Which properties of a planet can each method recover, and how does inclination
change the evidence?\\
Specialized plots &
Swept area, rotation curve, resonant angle, stellar tracks, gravitational-wave
strain, and conservation diagnostics &
What observation distinguishes a resonance from a near-integer ratio? Why does
a galaxy require more gravitating mass than its light implies?\\
Exports &
Screenshots, clips, report PDF, state link, and unit-labelled CSV &
Can another student or instructor reproduce the setup and inspect the evidence?\\
\hline
\end{tabular}
\end{table*}

Three observing panels are especially useful for connecting dynamics to
astronomical inference. A transit light curve, a radial-velocity curve, and an
astrometric track are generated from the same evolving system and the same
observer geometry. Moving the observer therefore changes all three consistently:
a planet can cease to transit, the radial-velocity amplitude changes with
inclination, and the astrometric line opens into an ellipse. Students encounter
selection effects and degeneracies as outcomes of a controlled experiment,
rather than only as warnings appended to a formula.

\begin{figure*}
\centering
\IfFileExists{figures/exoplanet-observing.png}
  {\includegraphics[width=0.96\textwidth]{figures/exoplanet-observing.png}}
  {\figureplaceholder{2.40in}{The Exoplanet Characterization Lab with the shared
  observer visible on the simulation canvas and the transit, radial-velocity,
  and astrometry panels open. Include plotted data and axis labels; avoid
  overlapping dialogs.}}
\caption{One simulated planetary system viewed through three observational
methods. Because the panels share the same observer, changing the viewing
geometry immediately exposes which inferences are robust and which depend on
inclination or alignment.}
\label{fig:exoplanets}
\end{figure*}

The application also makes numerical reliability visible when it is relevant.
Students may choose among three integrators and inspect changes in total energy
and angular momentum. Dedicated investigations distinguish a physical
sensitivity to initial conditions from ordinary phase drift or numerical error.
The purpose is not to teach numerical analysis in every activity; it is to let
students ask whether a surprising result survives an appropriate reliability
check.

\section{Three Paths Through the Simulation}
\label{sec:examples}

\subsection{Discovering Kepler's Laws}

In the Kepler investigation, students begin with the geometry of an ellipse.
They locate the star relative to the two foci, alter eccentricity, and compare
circular and eccentric orbits. The simulation then displays equal-area sectors
over equal time intervals while the student independently measures the fastest
and slowest portions of the orbit. Finally, periods and semimajor axes from
multiple systems are plotted to recover a slope of $3/2$ in log space.

This sequence connects picture, measurement, and equation. The second law is not
reduced to the statement that a planet moves faster when it is closer. The
learner sees the swept area, times the motion, and relates the result to angular
momentum. The third law is not supplied as an exponent to verify; it is the
relationship that emerges from the student's graph. The same content can be
used as a five-minute prediction and demonstration, a roughly twenty-minute
guided activity, or a full laboratory investigation.

\subsection{Becoming an Exoplanet Observer}

The exoplanet sequence begins with the shadow of HD~209458~b, a landmark
transiting system \citep{charbonneau2000}. Students measure a transit depth,
infer a radius, examine limb darkening, determine a period, and explore how
blended light from a companion star biases the result. The radial-velocity
investigation then uses the star's reflex motion to estimate the planet's mass.
Combining the two measurements produces a bulk density and a more physically
meaningful description of the planet.

Later investigations shift the question from parameter recovery to experimental
design. Students receive a finite number of observing nights, choose a cadence,
and compare schedules under the same signal and noise realization. They see how
aliasing, weather losses, baseline, and sampling affect whether a planet is
detected. The Goldilocks investigation then changes orbital distance, stellar
luminosity, and eccentricity to examine incident flux and the limits of the
phrase habitable zone, using established habitable-zone relations as the
scientific context \citep{kopparapu2013}.

Together, these activities reproduce the logic of observation: no single panel
reveals everything, geometry matters, and the observing strategy can determine
which conclusions are possible.

\subsection{When the Visible Matter Is Not Enough}

The Missing Mass investigation asks students to weigh systems twice. First they
sum the matter that is visible; then they infer gravitating mass from orbital
motion. These estimates agree in a Solar System-like case but diverge for a
galaxy and a galaxy cluster. Students arrange visible mass, inspect the rotation
curve it produces, and attempt to reproduce a measured flat curve. Adding an
extended halo changes the force law used by the orbiting tracers and allows the
curve to remain flat at large radius, connecting the experiment with modern
rotation-curve data sets such as SPARC \citep{lelli2016}.

This activity illustrates a broader goal of \Gravitas: students should
experience why an inference was made. The conclusion that galaxies contain dark
matter becomes the resolution of a failed mass model rather than a fact
presented without its observational pressure.

\section{Designed for Real Classrooms}
\label{sec:instructors}

An educational tool is only useful if it fits the time, equipment, and workflow
available to an instructor. \Gravitas\ supports several levels of use:

\begin{itemize}
\item \textbf{Live demonstration.} Lecture mode enlarges text, selects a
high-contrast daylight theme, supplies a spotlight pointer, and lets an
instructor step through a prepared sequence with arrow keys.
\item \textbf{Embedded figure.} A share link with embed mode produces a
chrome-free interactive view suitable for a course page or learning-management
system.
\item \textbf{Short activity.} Curated routes extract a coherent
prediction--experiment--explanation cycle from a longer investigation for use
during a single class meeting.
\item \textbf{Laboratory assignment.} A full investigation saves progress and
produces a student report that can be submitted through the course's existing
system.
\item \textbf{Open exploration or project.} The sandbox, A/B bench, share links,
and data export support questions designed by the learner or instructor.
\end{itemize}

For each guided investigation, the instructor collection includes learning
objectives, expected observations, common wrong turns, and an answer key
generated from the same lesson source used by students. An adopter's guide and
curriculum map help place activities into an existing course. Public teaching
pages explain the workflow and provide demonstration links before an instructor
requests access to protected answer materials.

\begin{figure*}
\centering
\IfFileExists{figures/breadth-and-classroom.png}
  {\includegraphics[width=0.96\textwidth]{figures/breadth-and-classroom.png}}
  {\figureplaceholder{2.35in}{A clean four-panel montage showing the scenario
  gallery plus three contrasting learning environments: the Missing Mass
  rotation-curve view, Listening to Spacetime or GW150914, and the stellar
  H--R diagram. Use labels (a)--(d) and retain readable axes.}}
\caption{\Gravitas\ extends beyond a conventional orbit sandbox. The final
montage should communicate both the searchable breadth of the scenario gallery
and representative activities in galactic dynamics, gravitational waves, and
stellar astrophysics.}
\label{fig:breadth}
\end{figure*}

Access considerations are treated as instructional infrastructure rather than
optional polish. The interface and all investigations are available in English
and Spanish. Keyboard navigation, screen-reader labels, reduced-motion
behavior, multiple themes, and automated accessibility checks support a wider
range of learners. The application can operate offline after its assets have
been cached, including on lower-powered devices with a reduced visual-quality
tier. Because all computation and progress storage occur locally, students are
not required to create an account or send their work to a third-party service.

These features also make assignments more resilient. A link can reproduce the
starting state across a large class. A screenshot documents its own time and
scale. A CSV names the units in its columns. A report gathers the student's
predictions, measurements, graphs, and explanations. The goal is to reduce the
amount of custom technical scaffolding an instructor must build before the
physics can begin.

\section{Scientific Scope, Transparency, and Validation}
\label{sec:scope}

\Gravitas\ is an educational model, not a professional research code. Its core
dynamics are Newtonian and two-dimensional. Point-mass gravity is integrated in
a plane; close encounters use softening; and some scenarios allow selected
bodies to feel gravity without exerting it so that large tracer populations
remain responsive. Collisions are represented as perfectly inelastic mergers.
Relativistic quantities such as Schwarzschild radius may be evaluated
analytically, but general-relativistic orbital dynamics are not solved.
Gravitational-wave inspiral uses an educational prescription, and visual
elements such as accretion disks, jets, and the optional rubber-sheet spacetime
view are illustrative. Body sizes are compressed for visibility and must not be
measured from their drawn diameters.

These boundaries are stated prominently on the public model page. Features are
labelled as simulated, analytic, approximated, illustrative, or not modeled.
That vocabulary is also pedagogically useful: students can ask whether a
particular conclusion follows from the evolving equations, a closed-form
relationship, or a visual analogy. Instructors can decide whether the model is
appropriate for the learning objective before assigning it.

A public validation report tests orbital periods, Kepler relationships,
conservation laws, convergence behavior, escape speed, exoplanet observables,
habitable-zone quantities, rotation curves, black-hole scaling relations, and
stored parameters for real systems against stated expectations or published
sources. Scenario-stability checks audit the shipped catalogue, while browser
tests exercise the student workflows. These checks cannot establish that the
software is error-free, nor can they substitute for classroom evaluation. They
do make the scientific claims, tolerances, known approximations, and failure
conditions inspectable.

Version \Version\ should therefore be evaluated along two separate dimensions.
Scientific verification asks whether the model produces the quantity it claims
to produce within an appropriate tolerance. Educational evaluation asks how
students and instructors use the environment and whether the intended
reasoning appears in their work. The first is included in the release process.
Formal evidence for the second is future work. The teaching site provides a
local feedback template so adopters can record course context, time on task,
friction points, and observed student reasoning without transmitting student
data.

\section{Availability and Reuse}
\label{sec:availability}

The simulation is available at \url{https://gravitas-sim.online}. Source code,
documentation, model limitations, validation materials, and issue tracking are
available at
\url{https://github.com/gravitas-sim/gravitas-sim.github.io}. The project is
released under the MIT License.

The archival record for version \Version\ will be deposited with Zenodo before
the final manuscript is submitted:

\begin{center}
\textbf{DOI:} \texttt{\ReleaseDOI}\\
\textbf{Zenodo record:} \texttt{\ZenodoRecord}
\end{center}

The placeholders above are intentionally explicit so they can be found and
replaced after the GitHub release has been archived. The tagged release,
Zenodo record, manuscript version, scenario and investigation counts, and final
screenshots should all describe the same \Version\ state.

Instructors may use the built-in investigations as written, assign selected
steps, or create their own activities from shared scenarios and exported data.
The repository welcomes reports from classroom use, translations, new
scenarios, and proposed investigations. The most valuable feedback is concrete:
what learners predicted, what they measured, where an interface choice
interrupted the reasoning, and what explanation emerged after the experiment.

\section{Conclusion}
\label{sec:conclusion}

\Gravitas\ gives students a laboratory for systems that ordinary laboratories
cannot hold. A planet can be moved, a viewpoint changed, a dark halo removed, a
binary replayed from the same initial state, or an observing schedule redesigned
within minutes. More importantly, the environment asks students to make those
changes for a reason: to test a prediction and construct an explanation from
measured evidence.

For an instructor, the entry cost is deliberately low: open a link, choose a
scenario or investigation, and use it as a projected demonstration, an embedded
interactive figure, a short activity, or a complete lab. For a student, the
invitation is equally direct: build a system, run it, measure it, and find out
whether the universe in the model behaves as expected. \Gravitas\ is available
now at \url{https://gravitas-sim.online}, and the versioned \Version\ archive
will be citable at the DOI above when the release is minted.

\begin{acknowledgments}
The author welcomes feedback from instructors and students who use \Gravitas\
in a course. Development and maintenance information will be added here before
submission if institutional or funding acknowledgment is required.
\end{acknowledgments}

\software{\Gravitas\ \citep{gravitas_zenodo}}

\bibliography{references}
\bibliographystyle{aasjournal}

\end{document}
